% =====================================================================
% Chapitre 3 — Tests unitaires (C2.2.2) + Sécurité/Accessibilité/
% Évolutivité (C2.2.3) — ÉLIMINATOIRES
% =====================================================================

\section{Tests unitaires}\label{sec:tests-unitaires}
\competence{C2.2.2}{Jeu de tests unitaires couvrant une fonctionnalité demandée}{OUI}

Cette section montre la mise en place d'un harnais de tests unitaires couvrant les fonctionnalités principales, permettant de détecter automatiquement les régressions.

\subsection{Frameworks de test retenus}
Le backend NestJS est testé avec Jest, intégré nativement via \texttt{@nestjs/testing}, tandis que le frontend Angular est testé avec Jasmine et Karma (outils natifs). Jest est le standard de l'écosystème Node.js et facilite le \emph{mocking} des dépendances injectées.

\subsection{Stratégie de test}
La logique métier des tournois est en grande partie extraite dans des fonctions utilitaires pures (calcul du classement, détermination du résultat d'un match, appariements, générateurs de format, statistiques)~: sans état ni dépendance injectée, elles sont testables directement, sans mock.

Le service d'orchestration \texttt{TournoisService} est quant à lui couvert par une série de suites ciblant chaque cas d'usage (démarrage, saisie, validation, terminaison, annulation, statistiques). Les points d'entrée sont également vérifiés~: le contrôleur (délégation au service et présence du guard d'authentification), le guard \texttt{GardeAuthentification} lui-même et la validation des DTO.

Cette stratégie se traduit par une couverture backend, mesurée par Jest (\texttt{npm run test:cov}) et publiée comme artefact du pipeline d'intégration continue (format Cobertura), de 82,81~\% des instructions (1479/1786) pour 208 tests répartis sur 29 suites, le détail par mesure figurant en annexe. L'effort ayant porté sur la logique métier sensible, le cœur du module tournois dépasse 90~\%, au-delà de la cible de 70~\% fixée au chapitre~1, la couverture du frontend Angular restant à compléter avant la mise en production. Le rapport complet, détaillé fichier par fichier, est fourni en \hyperref[chap:annexe-couverture]{annexe~\ref*{chap:annexe-couverture}}.

Le tableau suivant récapitule les unités testées et la logique couverte.

\begin{table}[H]
    \centering
    \small
    \renewcommand{\arraystretch}{1.6}
    \begin{tabular}{|>{\raggedright\arraybackslash}p{0.34\linewidth}|>{\raggedright\arraybackslash}p{0.56\linewidth}|}
        \hline
        \textbf{Unité testée} & \textbf{Logique couverte} \\
        \hline
        \texttt{TournoisService} & Orchestration du cycle de vie, répartie en plusieurs suites~: création, démarrage, saisie et validation des matchs, terminaison de phase, annulation et reprise, statistiques. \\
        \hline
        \texttt{classement.util} & Calcul du classement (points, victoires, défaites, différentiel) et départage des égalités. \\
        \hline
        \texttt{calculer-resultat-match} & Détermination du vainqueur d'un match à partir des scores de séries. \\
        \hline
        Générateurs de format et bracket & Construction des matchs et du bracket pour l'élimination directe, les poules (round-robin) et la ronde suisse. \\
        \hline
        \texttt{apparier} / \texttt{apparier-ronde-suisse} & Appariement des affrontements (tête de série, aléatoire, ronde suisse). \\
        \hline
        \texttt{statistiques.util} / \texttt{statistiques-equipe.util} & Calcul des statistiques du tournoi et par équipe. \\
        \hline
        Points d'entrée (\texttt{TournoisController}, \texttt{GardeAuthentification}, DTO) & Délégation du contrôleur au service et présence du guard, rejet 401 du guard en l'absence de session valide, et validation des données entrantes par les DTO. \\
        \hline
    \end{tabular}
\end{table}

Les mécanismes internes de better-auth (comptes, hachage des mots de passe, sessions) ne sont en revanche pas retestés~: ils relèvent de la librairie tierce, testée en amont par ses mainteneurs. Seule son intégration applicative, le guard \texttt{GardeAuthentification}, fait l'objet de tests de notre côté.

\subsection{Cas de test représentatif}
Le cas suivant, tiré de \texttt{classement.util.spec.ts} et construit selon le pattern Arrange / Act / Assert, illustre le départage des égalités dans \texttt{classement.util} (confrontation directe, puis différence de score, puis nombre de victoires). Quatre équipes sont à égalité de points par deux~: la différence favoriserait l'équipe~2, mais l'équipe~1 l'a battue en confrontation directe, critère prioritaire que vérifie l'assertion finale.

\begin{lstlisting}
// Arrange : 1 et 2 totalisent 6 points, 3 et 4 en totalisent 3.
// La difference de score favorise 2, mais 1 a battu 2 en direct.
const classement = calculerClassement(participants4, [
  matchTermine(1, 2, 3, 2),
  matchTermine(1, 3, 3, 2),
  matchTermine(4, 1, 5, 0),
  matchTermine(2, 3, 5, 0),
  matchTermine(2, 4, 5, 0),
  matchTermine(3, 4, 5, 0),
]);
// Act + Assert : la confrontation directe place 1 devant 2.
expect(classement.map((l) => l.participantId)).toEqual([1, 2, 3, 4]);
\end{lstlisting}

\subsection{Portée et limites de la couverture}
Un taux de couverture élevé ne constitue pas à lui seul une preuve de qualité. Une couverture importante peut coexister avec des défauts si les tests n'éprouvent pas la logique métier réellement sensible. L'effort a donc porté en priorité sur des tests orientés métier plutôt que sur la recherche d'un pourcentage. Chaque générateur de format est par exemple vérifié sur le nombre de matchs qu'il doit produire~: une phase de poules (round-robin) à dix participants doit générer neuf journées de cinq matchs, soit quarante-cinq matchs, et toute autre valeur fait échouer le test.

La valeur de cette approche s'est confirmée lors de la recette. Le commanditaire a rejoué de vrais tournois (championnats esport de type LFL), et le seul défaut métier rencontré concernait le classement de la ronde suisse, qui réutilisait à tort la logique de départage des poules. Ce cas avait échappé aux tests en raison d'une assertion mal écrite, ce qui illustre qu'un logiciel testé peut malgré tout comporter un défaut. Depuis, ce scénario est verrouillé par un test dédié (voir la \hyperref[sec:plan-correction]{partie~\ref*{sec:plan-correction}}), et toute régression future est détectée automatiquement, sans nouvelle sollicitation du commanditaire. La couverture prend ainsi tout son sens combinée à la recette~: les tests rejouent en continu un besoin métier déjà validé par le client.

\livrable{un jeu de tests unitaires couvrant une fonctionnalité demandée}

\sanssautdesection
\section{Sécurité, accessibilité et évolutivité}
\competence{C2.2.3}{Mesures de sécurité et actions d'accessibilité}{OUI}

Cette section présente les mesures de sécurité de l'application, les actions d'accessibilité engagées et l'évolutivité de son architecture. Le référentiel OWASP y sert de grille de lecture pour couvrir les principaux risques, sans que ces mesures aient encore fait l'objet d'un audit de sécurité ou d'un test d'intrusion externe.

\subsection{Sécurité — OWASP Top 10}
Chaque risque de l'OWASP Top 10 (2025) est associé à une mesure concrète mise en place dans BrackX. Ces bonnes pratiques applicatives réduisent la surface d'exposition sans prétendre à une conformité certifiée, qui supposerait un audit dédié.

\begin{longtable}{|>{\raggedright\arraybackslash}p{0.30\linewidth}|>{\raggedright\arraybackslash}p{0.60\linewidth}|}
    \hline
    \textbf{Risque (OWASP Top 10~: 2025)} & \textbf{Mesure mise en place} \\
    \hline
    \endhead
    A01 --- Broken Access Control & Toute route protégée passe par le guard \texttt{GardeAuthentification}, qui valide la session better-auth et rejette toute requête sans session valide (401). Chaque endpoint de tournoi vérifie en outre que la ressource appartient à l'utilisateur courant, ce qui empêche l'accès aux tournois d'autrui. Le rôle (\texttt{USER} / \texttt{ADMIN} / \texttt{DEV}) est exposé en lecture seule (\texttt{input: false})~: un client ne peut pas s'auto-attribuer un rôle privilégié. \\
    \hline
    A02 --- Security Misconfiguration & Variables d'environnement externalisées et non versionnées (\texttt{.env}). CORS configuré explicitement (origine limitée au frontend, méthodes et en-têtes restreints) et \texttt{trustedOrigins} better-auth aligné sur la même origine. \\
    \hline
    A03 --- Software Supply Chain Failures & Versions des dépendances figées dans \texttt{package-lock.json} et installées via \texttt{npm ci} (builds reproductibles). L'image de production n'est construite et publiée que par le pipeline CI/CD GitLab, à partir du \texttt{Dockerfile} versionné. Le Dependency Scanning GitLab est prévu pour automatiser la détection des dépendances vulnérables. \\
    \hline
    A04 --- Cryptographic Failures & Mots de passe hachés par better-auth et stockés dans la table \texttt{account}, jamais en clair. Les sessions sont signées avec un secret (\texttt{BETTER\_AUTH\_SECRET}) externalisé, et l'identifiant de session est transporté par un cookie \texttt{httpOnly} + \texttt{SameSite=Lax}, inaccessible au JavaScript. \\
    \hline
    A05 --- Injection & Prisma génère des requêtes paramétrées nativement, rendant les injections SQL impossibles. Le \texttt{ValidationPipe} global de NestJS (\texttt{whitelist}, \texttt{forbidNonWhitelisted}, \texttt{transform}) valide les DTO et rejette tout champ non attendu. \\
    \hline
    A06 --- Insecure Design & Architecture pensée selon le principe du moindre privilège~: le rôle par défaut est \texttt{USER} et ne peut être modifié par le client. Les endpoints d'administration sont restreints par le guard \texttt{GardeRoles} (\texttt{@Roles('ADMIN')}), qui renvoie un 403 si le rôle \texttt{ADMIN} fait défaut. \\
    \hline
    A07 --- Authentication Failures & Authentification déléguée à better-auth~: session limitée à 7 jours avec prolongement glissant, politique de mot de passe imposée (8 caractères minimum, une majuscule, une minuscule et un caractère spécial) vérifiée côté client et serveur, et réinitialisation par token à usage unique envoyé par email. \\
    \hline
    A08 --- Software or Data Integrity Failures & Seul le pipeline CI/CD déploie en production~: les images Docker sont versionnées par tag (\texttt{:vX.Y.Z}) et tirées d'un registre maîtrisé, sans mise à jour automatique depuis une source non vérifiée. \\
    \hline
    A09 --- Security Logging and Alerting Failures & La gestion des sessions et des échecs d'authentification est assurée par better-auth et remontée par NestJS. La centralisation des journaux, les alertes et la supervision (Grafana) sont prévues, un middleware de journalisation des erreurs 4xx/5xx restant à ajouter. \\
    \hline
    A10 --- Mishandling of Exceptional Conditions & Les erreurs sont traitées de manière contrôlée~: le \texttt{ValidationPipe} renvoie des réponses 400 structurées et le guard d'authentification un 401 explicite, sans exposer de détail interne. En production, la verbosité des messages est réduite. Un filtre d'exception NestJS centralisant le format des erreurs (sans fuite de \emph{stack trace}) reste à généraliser. \\
    \hline
\end{longtable}

L'extrait suivant illustre la première ligne de défense (A01)~: le guard \texttt{GardeAuthentification} interroge better-auth à partir des cookies de la requête et rejette tout accès dépourvu de session valide par un code 401, avant d'attacher l'utilisateur courant à la requête.

\begin{lstlisting}
@Injectable()
export class GardeAuthentification implements CanActivate {
  async canActivate(context: ExecutionContext): Promise<boolean> {
    const requete = context.switchToHttp().getRequest();
    // Recupere la session aupres de better-auth (cookie httpOnly)
    const session = await authentification.api.getSession({
      headers: fromNodeHeaders(requete.headers),
    });
    // Sans session valide : acces refuse (401)
    if (!session) {
      throw new UnauthorizedException('Non authentifie');
    }
    // Session valide : on attache l'utilisateur courant a la requete
    requete.user = session.user;
    requete.session = session.session;
    return true;
  }
}
\end{lstlisting}

\subsection{Données personnelles et conformité RGPD}
La seule donnée personnelle collectée est l'adresse email, qui sert d'identifiant de connexion~: aucune autre information nominative n'est demandée, ce qui limite l'exposition au regard du RGPD. Le volet conformité proprement dit reste toutefois à développer et constitue un chantier identifié avant la mise en production, et non une fonctionnalité déjà livrée. Deux mesures sont prévues~: l'anonymisation des tournois laissés inactifs, par exemple celle des équipes après un mois sans connexion du propriétaire, et la désactivation des comptes inactifs après deux à six mois sans connexion.

\subsection{Accessibilité — RGAA 4.1}
Le référentiel retenu est le RGAA 4.1, référence officielle française pour l'accessibilité numérique. Plusieurs mesures en rapprochent le prototype. Tous les éléments interactifs (boutons, liens, champs) portent un attribut \texttt{aria-label} ou \texttt{aria-labelledby} décrivant leur fonction, et les formulaires associent explicitement chaque label à son champ (\texttt{for}/\texttt{id}). Le contraste des textes vise le ratio minimum de 4.5:1 préconisé par le WCAG AA. La navigation au clavier est fonctionnelle sur l'ensemble de l'application, avec une tabulation logique et un focus visible, et les messages d'erreur de formulaire sont annoncés aux lecteurs d'écran via \texttt{aria-live}.

Ces éléments constituent un rapprochement du référentiel, non une conformité mesurée~: la navigation au clavier a été éprouvée lors des tests structurels de la recette, mais aucun audit RGAA formel ni test avec lecteur d'écran n'a encore été conduit. Un audit outillé (Lighthouse, axe-core) reste à mener pour chiffrer le niveau réellement atteint. Une limite est déjà identifiée~: la vue bracket en arbre est difficile à rendre pleinement accessible, ce pour quoi une alternative textuelle sous forme de tableau de résultats est proposée en complément.

\subsection{Évolutivité}
L'architecture modulaire permet d'ajouter des fonctionnalités sans refonte, ce que le développement a confirmé~: les statistiques ont été ajoutées à partir des données déjà structurées en base, via de nouveaux endpoints et des fonctions utilitaires dédiées, sans modifier le cœur existant.

Cette évolutivité tient surtout à la conception du cœur métier. Les trois formats disponibles (élimination directe, poules et ronde suisse) couvrent la grande majorité des tournois. Les cas restants s'ajoutent de façon incrémentale, car le back-end sait déjà créer les matchs, propager les qualifiés et calculer les classements~: un match de départage à égalité, une petite finale ou un tableau des perdants (lower bracket) ne demandent qu'un match supplémentaire et sa visualisation, sans toucher à la logique existante. Le coût principal de ces ajouts porte donc sur l'affichage côté front, non sur le back-end.

Le point le plus structurant est que le schéma de données porte déjà les liaisons entre matchs~: le vainqueur d'un match alimente automatiquement le match suivant. Aujourd'hui, un template génère ces liaisons selon des règles prédéfinies, mais rien n'empêche de les définir librement, par exemple par glisser-déposer. Une telle édition libre ne remettrait pas en cause le back-end, dont la logique de propagation resterait identique~: l'effort se concentrerait sur l'expérience de construction côté front. Les évolutions planifiées suivent ce principe~: le partage public d'un tournoi ne demande qu'un champ \texttt{isPublic} et un guard de lecture non authentifiée, et les notifications temps réel s'appuieront sur les WebSockets natifs de NestJS, sans impacter les modules existants.

L'évolutivité fonctionnelle ne doit pas être confondue avec la montée en charge. L'architecture actuelle regroupe le frontend, le back-end et la base dans une même image, ce qui convient au volume attendu mais n'est pas dimensionné pour une forte charge simultanée, dont le point de rupture serait le couple back-end et base de données. Une montée à grande échelle passerait par la segmentation des services et une orchestration (Kubernetes) dimensionnant chaque service selon des seuils de charge. Ce chantier, surdimensionné pour le nombre d'utilisateurs visé, est écarté à ce stade et relève de la supervision étudiée au Bloc 4.

\livrable{la présentation des mesures de sécurité et des actions d'accessibilité aux personnes en situation de handicap}
